\chapter{The analytic criterion}\label{ch:an}

This chapter covers the analytic end of the archive's chain. A one-sided
bound $G(x)\le x^{\theta}$ for a positive combination $G$ of smoothed prime
sums, assumed only on windows $[Y,Y^{\lambda}]$ for arbitrarily large $Y$,
forces the corresponding Dirichlet $L$-functions to have no zero in
$\Re s>\theta$ (a windowed form of Landau's theorem on Dirichlet integrals
with nonnegative coefficients). The archive proves this windowed theorem
in full in \arch{Development/AutomorphicArithmetic/DedekindLFunctions/LandauTheorems/}.
It then combines it with two inputs from other parts of the archive, a
Chebyshev-type inequality for the combination on the scales of a family of
host data and the existence of such hosts, to obtain zero-freeness in
$\Re s>1/2$ for every primitive Dirichlet character
(\texttt{ZeroFreeAll}).

\section{Smoothed prime sums}

\begin{definition}[Primitive characters]\label{an:pchar}
\lean{Formalism.HC.PChar, Formalism.HC.PChar.one, Formalism.HC.PChar.inv, Formalism.HC.PChar.of, Formalism.HC.PChar.IsReal}
A \emph{primitive character} is a pair $(N,\chi)$ with $N\ge1$ and $\chi$ a
primitive Dirichlet character modulo $N$ with values in $\CC$; two such
pairs are equal only if the moduli and the characters agree. The trivial
primitive character $\mathbf 1$ is the unique character modulo $1$. The
inverse of $(N,\chi)$ is $(N,\chi^{-1})$. For any Dirichlet character
$\chi$ modulo $N$, $\operatorname{of}(\chi)$ is the primitive character
inducing $\chi$, of modulus the conductor of $\chi$. A primitive character
is \emph{real} if all its values are real.
\end{definition}

\begin{definition}[Smoothed Chebyshev functions, real form]\label{an:smoothed}
\lean{Formalism.Sz, Formalism.Gz, Formalism.Schi, Formalism.Gchi}
Let $\Lambda$ be the von Mangoldt function and $\gamma$ Euler's constant.
For $x\in\R$ and a Dirichlet character $\chi$ modulo $N\ge1$ put
\begin{align*}
\Sz(x)&=\sum_{1\le n\le x}\Lambda(n)\Bigl(\frac xn-1\Bigr), &
\Gz(x)&=\Sz(x)-x\log x+(1+\gamma)x,\\
S_{\chi}(x)&=\sum_{1\le n\le x}\Lambda(n)\,\Re\chi(n)\Bigl(\frac xn-1\Bigr), &
G^{\R}_{\chi}(x)&=S_{\chi}(x)+x\,\Re\frac{L'(1,\chi)}{L(1,\chi)},
\end{align*}
where $L(s,\chi)$ is Mathlib's Dirichlet $L$-function (so $\log$ and
$L(1,\chi)$ take Mathlib's values at points where the classical object is
undefined).
\end{definition}

\begin{definition}[Smoothed Chebyshev functions, complex form]\label{an:gchi}
\lean{Formalism.HC.An.logDerivativeAtOne, Formalism.HC.An.G, Formalism.HC.An.sumLamChi, Formalism.HC.An.gChi}
For $y\in\R$ and a Dirichlet character $\chi$ modulo $N\ge1$ put
$\ell(\chi)=L'(1,\chi)/L(1,\chi)$ and
\[
G(y)=\sum_{1\le n\le y}\Lambda(n)\max\Bigl(\frac yn-1,0\Bigr)-y(\log y-1-\gamma),
\qquad
\Sigma_{\chi}(y)=\sum_{1\le n\le y}\Lambda(n)\chi(n)\max\Bigl(\frac yn-1,0\Bigr).
\]
Define $g_{\chi}(y)=G(y)$ if $\chi$ is the trivial character of its modulus,
and $g_{\chi}(y)=\Sigma_{\chi}(y)+y\,\ell(\chi)$ otherwise.
\end{definition}

\begin{lemma}\label{an:gz-eq}
\lean{UpperTheory.Comparison₂.Gz_eq}
\uses{an:smoothed, an:gchi, an:charfamily}
For every real $y$, $\Gz(y)=G(y)$, and $G(y)$ equals the function $G$ of
the character-family development (Definition~\ref{an:charfamily}).
\end{lemma}
\begin{proof}
\uses{an:smoothed, an:gchi}
For $1\le n\le\lfloor y\rfloor$ one has $n\le y$, so
$\max(y/n-1,0)=y/n-1$; the two expressions then differ by rearranging the
linear terms. The sum over $0\le n\le\lfloor y\rfloor$ used in the
character-family version differs only by the term $n=0$, where
$\Lambda(0)=0$.
\end{proof}

\begin{lemma}\label{an:re-gchi}
\lean{Formalism.HC.Landau.re_gChi}
\uses{an:smoothed, an:gchi}
If $\chi\ne1$, then $\Re g_{\chi}(y)=G^{\R}_{\chi}(y)$ for every real $y$.
\end{lemma}
\begin{proof}
\uses{an:smoothed, an:gchi}
Termwise: $\Lambda(n)$ and $\max(y/n-1,0)=y/n-1$ are real for
$1\le n\le y$, and $\Re(y\,\ell(\chi))=y\Re\ell(\chi)$.
\end{proof}

\section{Levels and the weighted combination}

\begin{definition}[Levels]\label{an:level}
\lean{Formalism.HC.KSet, Formalism.HC.Level, Formalism.HC.Level.Nt, Formalism.HC.Level.wt, Formalism.HC.Level.Gm}
\uses{an:pchar, an:smoothed, an:gchi}
A \emph{level} $\Lambda$ consists of
\begin{itemize}
\item a finite set $\Kept(\Lambda)$ of primitive characters containing
$\mathbf 1$ and closed under $\psi\mapsto\psi^{-1}$;
\item an optional character $\psi_{K}$, which, when present, lies in
$\Kept(\Lambda)$, is odd and is real;
\item a natural number $o(\Lambda)$.
\end{itemize}
Its weights are $\operatorname{wt}(\psi)=1$ if $\psi=\psi_K$ and
$\operatorname{wt}(\psi)=(3o(\Lambda)+1)^{-1}$ otherwise, and its
combination is
\[
G_{\Lambda}(y)=\Gz(y)+\sum_{\psi\in\Kept(\Lambda),\ \psi\ne\mathbf 1}
\operatorname{wt}(\psi)\,\Re g_{\psi}(y).
\]
All weights are positive.
\end{definition}

\begin{definition}[Cyclic levels]\label{an:level-cyc}
\lean{Formalism.HC.Level.ofSub, Formalism.HC.Level.cyc, Formalism.HC.Level.mem_cyc}
\uses{an:level}
For a subgroup $\Xi$ of the Dirichlet characters modulo $M\ge1$, the level
$\Lambda_{\Xi}$ has $\Kept(\Lambda_{\Xi})=\{\operatorname{of}(\chi):\chi\in\Xi\}$,
no exceptional character, and $o(\Lambda_{\Xi})=|\Xi|$. For a primitive
character $\psi$ modulo $N$, the \emph{cyclic level} $\Lambda_{\psi}$ is
$\Lambda_{\Xi}$ for $\Xi$ the cyclic group generated by $\psi$; so its kept
set consists of the primitive characters inducing the powers of $\psi$,
$o(\Lambda_{\psi})$ is the order of $\psi$, every nontrivial weight equals
$(3o(\Lambda_{\psi})+1)^{-1}$, and $\psi\in\Kept(\Lambda_{\psi})$.
\end{definition}

\begin{definition}[Zero-freeness predicates]\label{an:zerofree}
\lean{Formalism.HC.PChar.ZeroFreeAbove, Formalism.HC.ZeroFreeOn, Formalism.HC.ZeroFreeAll, Formalism.HC.Level.ZeroFreeAbove, Formalism.HC.Level.ZeroFreeLevel}
\uses{an:pchar, an:level}
A primitive character $\psi$ is \emph{zero-free above} $\theta$ if
$L(s,\psi)\ne0$ for every $s\in\CC$ with $\Re s>\theta$. For a property $Q$
of primitive characters, $\mathrm{ZeroFreeOn}(Q)$ says that every $\psi$
with $Q(\psi)$ is zero-free above $1/2$, and $\mathrm{ZeroFreeAll}$ is
$\mathrm{ZeroFreeOn}$ of the trivially true property. A level is
zero-free above $\theta$ if every kept character is, and
$\mathrm{ZeroFreeLevel}(\Lambda)$ means zero-free above $1/2$.
\end{definition}

\begin{remark}[the point $s=1$]\label{an:rem-junk}
For $\psi=\mathbf 1$, $L(s,\mathbf 1)$ is Mathlib's $\zeta(s)$, and the
predicates of Definition~\ref{an:zerofree} include the point $s=1$. There
the statement $\zeta(1)\ne0$ is about Mathlib's assigned value at the pole,
which is nonzero. Inside the Landau argument this point is handled by
Mathlib's nonvanishing of $\zeta$ on $\Re s\ge1$, which covers $s=1$
(Lemma~\ref{an:nonvanish-re1}).
\end{remark}

\begin{lemma}\label{an:zerofree-levels}
\lean{Formalism.HC.zeroFreeOn_of_levels, Formalism.HC.zeroFreeOn_of_cyc}
\uses{an:zerofree, an:level-cyc}
Let $Q$ be a property of primitive characters. If every $\psi$ with
$Q(\psi)$ lies in the kept set of some level $\Lambda$ with
$\mathrm{ZeroFreeLevel}(\Lambda)$, then $\mathrm{ZeroFreeOn}(Q)$. In
particular this holds if $\mathrm{ZeroFreeLevel}(\Lambda_{\psi})$ for every
$\psi$ with $Q(\psi)$.
\end{lemma}
\begin{proof}
\uses{an:zerofree, an:level-cyc}
Immediate from the definitions; for the second statement use
$\psi\in\Kept(\Lambda_{\psi})$.
\end{proof}

\section{Windows and the window criterion}

\begin{definition}[Windows]\label{an:windows}
\lean{Formalism.HC.HasWindows}
A set $S\subset\R$ \emph{has windows of exponent} $\lambda$ if for every
$y\in\R$ there is $Y\ge y$ such that every $x$ with $Y\le x\le Y^{\lambda}$
lies in $S$.
\end{definition}

\begin{lemma}\label{an:windows-mono}
\lean{Formalism.HC.HasWindows.mono, Formalism.HC.hasWindows_of_large}
\uses{an:windows}
Let $S$ have windows of exponent $\lambda$, and suppose every $x\ge x_0$
in $S$ lies in $S'$. Then $S'$ has windows of exponent $\lambda$. A set
containing $[x_0,\infty)$ has windows of every exponent.
\end{lemma}
\begin{proof}
\uses{an:windows}
Apply the hypothesis at $\max(y,x_0)$.
\end{proof}

\begin{definition}[The window criterion]\label{an:formb}
\lean{Formalism.HC.Level.FormBAt, Formalism.HC.Level.FormB}
\uses{an:level, an:windows, an:zerofree}
A level $\Lambda$ \emph{satisfies the window criterion at} $\lambda$
($\mathrm{FormBAt}$) if for every $\theta\in(0,1)$: if
$\{x: G_{\Lambda}(x)\le x^{\theta}\}$ has windows of exponent $\lambda$, then
$\Lambda$ is zero-free above $\theta$. It satisfies the window criterion
($\mathrm{FormB}$) if this holds for every $\lambda>1$.
\end{definition}

\begin{lemma}\label{an:formb-mono}
\lean{Formalism.HC.Level.FormBAt.mono}
\uses{an:formb}
If $\Lambda$ satisfies the window criterion at $\lambda$ and
$\lambda\le\lambda'$, it satisfies it at $\lambda'$.
\end{lemma}
\begin{proof}
\uses{an:formb, an:windows}
A window $[Y,Y^{\lambda'}]$ with $Y\ge1$ contains $[Y,Y^{\lambda}]$.
\end{proof}

\section{The windowed Landau theorem}

\begin{definition}[Character families]\label{an:charfamily}
\lean{Deduction.ramp, Deduction.Gchi, Deduction.G, Deduction.CharFamily, Deduction.Gm, Deduction.ZeroFree}
A \emph{character family} $F$ consists of a finite index set $I$, for
each $i\in I$ a Dirichlet character $\chi_i$ modulo $N_i\ge1$ with
$N_i=1$ or $\chi_i\ne1$, and weights $a_i>0$. For a Dirichlet character
$\chi$, the archive's second copy of the smoothed sum is
$\mathcal G_{\chi}(y)=\sum_{0\le n\le y}\Lambda(n)\chi(n)\max(y/n-1,0)-y\log
y+(1+\gamma)y$ if $\chi=1$, and the same sum plus $y\,L'(1,\chi)/L(1,\chi)$
otherwise. The family's combination is
$G_F(y)=\sum_i a_i\Re\mathcal G_{\chi_i}(y)$, and $F$ is \emph{zero-free
above} $\theta$ if $L(s,\chi_i)\ne0$ for all $i$ and all $\Re s>\theta$.
\end{definition}

\begin{remark}
The admissibility condition $N_i=1$ or $\chi_i\ne1$ excludes principal
characters of modulus larger than $1$, whose $L$-functions have a pole at
$s=1$ and a Mellin transform of a different shape.
\end{remark}

\begin{definition}[Window hypotheses]\label{an:windowhyp}
\lean{Deduction.WindowHyp, Deduction.gm, Deduction.WindowHypU}
\uses{an:charfamily}
Let $F$ be a character family and $\theta,\delta\in\R$. The \emph{window
hypothesis} $\mathrm{WH}(F,\theta,\delta)$ says: for every $\eta>0$ and
every $Y_0$ there are $Y\ge Y_0$ and a function $R$ on
$[Y,Y^{1+\delta}]$ with $R\ge0$, $G_F\le R$ there, $R(y)y^{-\theta-1}$
integrable on $[Y,Y^{1+\delta}]$ and
$\int_Y^{Y^{1+\delta}}R(y)y^{-\theta-1}\,dy\le Y^{\eta}$.
Its logarithmic form $\mathrm{WH}_U(g,\theta,\delta)$, for a function $g$
on $\R$, says: for every $\eta>0$ and $U_0$ there are $U\ge U_0$ and $R$
with $0\le R$, $g\le R$ on $[U,(1+\delta)U]$, $R(u)e^{-\theta u}$
integrable there and $\int_U^{(1+\delta)U}R(u)e^{-\theta u}\,du\le
e^{\eta U}$. Write $g_F(u)=G_F(e^u)$.
\end{definition}

\begin{lemma}[From a power bound to the window hypothesis]\label{an:windowhyp-rpow}
\lean{Deduction.windowHyp_of_rpow, Deduction.windowHyp_of_rpow_eventually}
\uses{an:windowhyp}
Let $\delta>0$ and $C\in\R$, and suppose that for every $Y_0$ there is
$Y\ge Y_0$ with $G_F(y)\le Cy^{\theta}$ for all $y\in[Y,Y^{1+\delta}]$.
Then $\mathrm{WH}(F,\theta,\delta)$ holds.
\end{lemma}
\begin{proof}
\uses{an:windowhyp}
Take $R(y)=|C|y^{\theta}$. Then $\int_Y^{Y^{1+\delta}}R(y)y^{-\theta-1}dy
=|C|\delta\log Y$, and $|C|\delta\log Y\le Y^{\eta}$ once $Y$ is large,
because $\log Y\le Y^{\eta/2}/(\eta/2)$. Apply the hypothesis at a large
enough $Y_0$.
\end{proof}

\begin{lemma}\label{an:windowhypu}
\lean{Deduction.windowHypU_of_windowHyp}
\uses{an:windowhyp}
If $\delta>0$ and $\mathrm{WH}(F,\theta,\delta)$, then
$\mathrm{WH}_U(g_F,\theta,\delta)$.
\end{lemma}
\begin{proof}
\uses{an:windowhyp}
Substitute $y=e^u$: the window $[Y,Y^{1+\delta}]$ becomes
$[\log Y,(1+\delta)\log Y]$ and $R(y)y^{-\theta-1}dy$ becomes
$R(e^u)e^{-\theta u}du$; $Y^{\eta}=e^{\eta\log Y}$.
\end{proof}

\begin{definition}[Pole data]\label{an:poledata}
\lean{Deduction.zeroSet, Deduction.PoleData}
\uses{an:charfamily}
For a character family $F$, let $Z(F)$ be the set of $\rho$ with
$\Re\rho>0$ such that $L(\rho,\chi_i)=0$ or $L(\bar\rho,\chi_i)=0$ for
some $i$. For $g:\R\to\R$, $\Phi:\CC\to\CC$ and $Z\subset\CC$, the
statement $\mathrm{PoleData}(g,\Phi,Z)$ consists of:
\begin{enumerate}
\item $g$ is measurable and $|g(u)|\le Ce^{2u}$ for $u\ge0$, for some $C$;
\item $\Phi(s)=\int_0^{\infty} g(u)e^{-su}\,du$ for $\Re s>2$;
\item every $\rho\in Z$ has $\Re\rho>0$, $\Phi$ is holomorphic on
$\{\Re s>0\}\setminus Z$, and $Z$ has no accumulation point in
$\{\Re s>0\}$ (every such point has a punctured neighbourhood disjoint
from $Z$);
\item at each $\rho\in Z$, $\Phi(s)=c/(s-\rho)+h(s)$ near $\rho$ with
$h$ holomorphic at $\rho$ and $c\ne0$, and if $\rho$ is real then $c$ is
real and positive.
\end{enumerate}
\end{definition}

\begin{lemma}[Growth]\label{an:gchi-growth}
\lean{Deduction.Gchi_growth}
\uses{an:charfamily}
For every Dirichlet character $\chi$ there is $C$ with
$|\mathcal G_{\chi}(y)|\le Cy^2$ for $y\ge1$.
\end{lemma}
\begin{proof}
\uses{an:charfamily}
The prime sum is at most $2y^2$ in absolute value, since $|\chi(n)|\le1$;
the remaining terms are $O(y\log y)$.
\end{proof}

\begin{lemma}[Conjugation]\label{an:gchi-conj}
\lean{Deduction.Gchi_conj, Deduction.LFunction_conj}
\uses{an:charfamily}
For a Dirichlet character $\chi$ and $s\in\CC$,
$L(\bar s,\chi^{-1})=\overline{L(s,\chi)}$, and
$\overline{\mathcal G_{\chi}(y)}=\mathcal G_{\chi^{-1}}(y)$ for real $y$.
\end{lemma}
\begin{proof}
\uses{an:charfamily}
The first identity is proved through the Hurwitz zeta functions, which
commute with complex conjugation; the second follows termwise, using
$\overline{\chi(n)}=\chi^{-1}(n)$ and the first identity for the
derivative at $1$.
\end{proof}

\begin{lemma}[Mellin transforms]\label{an:mellin}
\lean{Deduction.mellin_Gchi_of_ne_one, Deduction.mellin_Gchi_modOne, Deduction.poleData_change}
\uses{an:charfamily}
Let $\Re s>1$. If $\chi\ne1$, then
\[
\int_1^{\infty}\mathcal G_{\chi}(y)y^{-s-1}\,dy
=-\frac{L'(s,\chi)/L(s,\chi)-s\,L'(1,\chi)/L(1,\chi)}{s(s-1)}.
\]
For the character modulo $1$,
\[
\int_1^{\infty}\mathcal G_{\chi}(y)y^{-s-1}\,dy
=-\frac{\zeta'(s)/\zeta(s)+1/(s-1)+1-(1+\gamma)s}{s(s-1)}.
\]
Moreover $\int_0^{\infty} H(e^u)e^{-su}\,du=\int_1^{\infty} H(y)y^{-s-1}\,dy$
for every function $H$.
\end{lemma}
\begin{proof}
\uses{an:charfamily}
Termwise integration: $\int_1^{\infty}\max(y/n-1,0)y^{-s-1}dy
=n^{-s}/(s(s-1))$, the series $\sum\Lambda(n)\chi(n)n^{-s}$ is
$-L'(s,\chi)/L(s,\chi)$, and the elementary terms give
$\int_1^{\infty}\log y\cdot y^{-s}dy=(s-1)^{-2}$ and similar integrals. The last
identity is the substitution $y=e^u$.
\end{proof}

\begin{lemma}[Pole structure of one transform]\label{an:polechar}
\lean{Deduction.poleData_abstract, Deduction.poleData_char}
\uses{an:mellin, an:gchi-growth}
Let $\chi$ be a Dirichlet character modulo $N$ with $N=1$ or $\chi\ne1$.
There is $\Phi_{\chi}:\CC\to\CC$ such that
\begin{enumerate}
\item $\Phi_{\chi}(s)=\int_0^{\infty}\mathcal G_{\chi}(e^u)e^{-su}\,du$ for $\Re s>2$;
\item $\Phi_{\chi}$ is holomorphic at every $\rho\ne0$ with
$L(\rho,\chi)\ne0$;
\item at every $\rho\ne0$ there are $m\in\N$, with $m=0$ exactly when
$L(\rho,\chi)\ne0$, and $h$ holomorphic at $\rho$, such that
$\Phi_{\chi}(s)=m\cdot\bigl(-1/(\rho(\rho-1))\bigr)/(s-\rho)+h(s)$ near $\rho$;
\item every point has a punctured neighbourhood on which $L(\cdot,\chi)$
does not vanish.
\end{enumerate}
\end{lemma}
\begin{proof}
\uses{an:mellin, an:gchi-growth}
An abstract lemma treats an entire $g$ with $g(1)\ne0$ and a constant
$k$: the function
$-\bigl(g'/g(s)-s\,g'(1)/g(1)+k(1-s)\bigr)/(s(s-1))$ has a removable
singularity at $s=1$, is holomorphic where $g\ne0$ and $s\ne0$, and at a
zero of $g$ of order $m$ has the stated simple pole, because $g'/g$ has
residue $m$ there. For $\chi\ne1$ take $g=L(\cdot,\chi)$ (entire, with
$L(1,\chi)\ne0$) and $k=0$. For the character modulo $1$ take
$g(s)=(s-1)\zeta(s)$, extended by $g(1)=1$, and $k=1$; then $g'(1)=\gamma$,
and $g(w)\ne0\iff\zeta(w)\ne0$ for all $w$ (at $w=1$ by Mathlib's
$\zeta(1)\ne0$). Lemma~\ref{an:mellin} identifies the function with the
Laplace transform on $\Re s>2$.
\end{proof}

\begin{lemma}[Pole data of a family]\label{an:poledata-lemma}
\lean{Deduction.poleData}
\uses{an:poledata, an:polechar, an:gchi-growth, an:gchi-conj, an:nonvanish-re1}
For every character family $F$ there is $\Phi$ with
$\mathrm{PoleData}(g_F,\Phi,Z(F))$.
\end{lemma}
\begin{proof}
\uses{an:polechar, an:gchi-growth, an:gchi-conj, an:nonvanish-re1}
Since $\Re z=(z+\bar z)/2$ and $\overline{\mathcal G_{\chi}}=\mathcal
G_{\chi^{-1}}$, take $\Phi=\sum_i\frac{a_i}2(\Phi_{\chi_i}+\Phi_{\chi_i^{-1}})$.
Growth follows from Lemma~\ref{an:gchi-growth}, and the Laplace formula
from Lemma~\ref{an:polechar}(1). By Lemma~\ref{an:gchi-conj},
$L(\rho,\chi_i^{-1})=0$ iff $L(\bar\rho,\chi_i)=0$, so off $Z(F)$ every
summand is holomorphic. At $\rho\in Z(F)$ the residue is
$c=S\cdot(-1/(\rho(\rho-1)))$ with
$S=\sum_i\frac{a_i}{2}(m_{1,i}+m_{2,i})$, where $m_{1,i},m_{2,i}$ are the
orders of $L(\cdot,\chi_i)$ and $L(\cdot,\chi_i^{-1})$ at $\rho$; $S>0$
because some order is positive and all $a_i>0$. Also $\rho\ne1$ and, for
real $\rho$, $\rho<1$ by Lemma~\ref{an:nonvanish-re1}; then
$\rho(\rho-1)<0$, so $c$ is real and positive.
\end{proof}

\begin{lemma}[Nonvanishing on $\Re s\ge1$]\label{an:nonvanish-re1}
\lean{Deduction.charFamily_LFunction_ne_zero_of_one_le_re}
\uses{an:charfamily}
For a character family $F$, every $i$ and every $s$ with $\Re s\ge1$,
$L(s,\chi_i)\ne0$.
\end{lemma}
\begin{proof}
\uses{an:charfamily}
If $N_i=1$, $L(\cdot,\chi_i)=\zeta$ and Mathlib's nonvanishing of $\zeta$ on
$\Re s\ge1$ applies (it includes Mathlib's value at $s=1$). Otherwise
$\chi_i\ne1$ and Mathlib's theorem for nontrivial characters applies.
\end{proof}

\begin{lemma}[Taylor coefficients]\label{an:coeff}
\lean{Deduction.landau_coeff, Deduction.iteratedDeriv_laplace}
\uses{an:poledata}
Assume $\mathrm{PoleData}(g,\Phi,Z)$ and $\Re s>2$. Then for every
$k\in\N$,
\[
c_k(s):=\frac{(-1)^k\Phi^{(k)}(s)}{k!}=\int_0^{\infty} g(u)\frac{u^k}{k!}e^{-su}\,du.
\]
\end{lemma}
\begin{proof}
\uses{an:poledata}
Differentiate the Laplace integral under the integral sign, using
$|g(u)|\le Ce^{2u}$ for domination; $\Phi$ agrees with the integral on
the open half-plane $\Re s>2$.
\end{proof}

\begin{lemma}[Concentration of the Gamma kernel]\label{an:tails}
\lean{Deduction.landau_tails}
Let $\delta>0$. There are $\sigma>2$ and $C$ such that for every
$U\ge1$ some integer $k\ge U$ satisfies
\[
\int_{(0,\infty)\setminus[U,(1+\delta)U]}\frac{u^k}{k!}e^{(2-\sigma)u}\,du\le\frac{C}{\sigma^k}.
\]
\end{lemma}
\begin{proof}
Put $l=2+6/\delta$, $\sigma=l(1+\delta)(2+l)$ and
$k=\lceil(1+\delta/2)\sigma U\rceil$, so $k\ge U$. Bound the two tails
of the kernel by exponential tilting with parameter $l$ (a Chernoff
bound: $\mathbf 1_{u<U}\le e^{l(U-u)}$ and
$\mathbf 1_{u>V}\le e^{l(u-V)}$), evaluate the resulting Gamma integrals
exactly, and check that both are at most a constant times $\sigma^{-k}$
for this choice of $k$.
\end{proof}

\begin{lemma}[One-sided window inequality]\label{an:window-ineq}
\lean{Deduction.landau_window_ineq, Deduction.landau_kernel_max}
Let $g$ be measurable with $|g(u)|\le C_0e^{2u}$ for $u\ge0$, let
$\sigma>2$, $\theta<\sigma$, $0<U$, and let $R$ satisfy $0\le R$, $g\le R$
on $[U,V]$, with $R(u)e^{-\theta u}$ integrable on $[U,V]$. Then for all
$k\in\N$ and $\gamma\in\R$,
\begin{multline*}
\Bigl|\int_0^{\infty} g(u)\frac{u^k}{k!}e^{-(\sigma+i\gamma)u}du\Bigr|
+\Re\int_0^{\infty} g(u)\frac{u^k}{k!}e^{-\sigma u}du\\
\le\frac{2}{(\sigma-\theta)^k}\int_U^VR(u)e^{-\theta u}du
+2C_0\int_{(0,\infty)\setminus[U,V]}\frac{u^k}{k!}e^{(2-\sigma)u}du.
\end{multline*}
\end{lemma}
\begin{proof}
On $[U,V]$, $g\le R$ and $R\ge0$ give $|g|\le 2R-g$, and
$u^ke^{-\sigma u}/k!\le(\sigma-\theta)^{-k}e^{-\theta u}$ because
$t^k/k!\le e^t$ with $t=(\sigma-\theta)u$. Outside $[U,V]$ use
$|g|\le C_0e^{2u}$ for both integrals.
\end{proof}

\begin{lemma}[Coefficient bound from a window]\label{an:window-key}
\lean{Deduction.landau_window_key}
\uses{an:poledata, an:coeff, an:tails, an:window-ineq, an:windowhyp}
Assume $\mathrm{PoleData}(g,\Phi,Z)$, $0<\theta<1$, $\delta>0$ and
$\mathrm{WH}_U(g,\theta,\delta)$. Then there are $\sigma>2$ and $C$ such
that for every $\eta>0$ and $U_0$ there are $U\ge U_0$ and an integer
$k\ge U$ with
\[
|c_k(\sigma+i\gamma)|+\Re c_k(\sigma)\le\frac{Ce^{\eta U}}{(\sigma-\theta)^k}
\qquad\text{for all }\gamma\in\R.
\]
\end{lemma}
\begin{proof}
\uses{an:coeff, an:tails, an:window-ineq}
Take $\sigma,C$ from Lemma~\ref{an:tails}. Given $\eta,U_0$, the window
hypothesis gives a window $[U,(1+\delta)U]$ with $U\ge\max(U_0,1)$ and a
majorant $R$ with $\int R e^{-\theta u}\le e^{\eta U}$;
Lemma~\ref{an:tails} gives $k\ge U$ for that $U$. Insert both into
Lemma~\ref{an:window-ineq}, express the integrals as $c_k$ by
Lemma~\ref{an:coeff}, and bound $C/\sigma^k\le|C|e^{\eta U}/(\sigma-\theta)^k$.
\end{proof}

\begin{lemma}\label{an:pole-re-le}
\lean{Deduction.landau_pole_re_le, Deduction.landau_pole_re_le_bound}
\uses{an:poledata}
If $\mathrm{PoleData}(g,\Phi,Z)$, then $\Re\rho\le2$ for every $\rho\in Z$.
\end{lemma}
\begin{proof}
\uses{an:poledata}
For $\sigma>2$, $\Phi$ is bounded on $\Re s\ge\sigma$ by the Laplace
integral, while near a pole with nonzero residue it is unbounded.
\end{proof}

\begin{lemma}[Dichotomy]\label{an:dichotomy}
\lean{Deduction.landau_dichotomy, Deduction.landau_finite}
Let $Z\subset\CC$ have no accumulation point in $\Re s>0$, let
$\Re\rho<\sigma$ for all $\rho\in Z$, let $0<\theta<\sigma$, and let
$\rho_0\in Z$ with $\Re\rho_0>\theta$. Then either
\begin{enumerate}
\item there is $D\le\sigma-\theta$ such that the open disc $B(\sigma,D)$
contains no point of $Z$ but some $\rho_1\in Z$ has $\sigma-\Re\rho_1<D$;
or
\item there are a real $\beta\in Z$ with $\beta>\theta$ and $r$ with
$\sigma-\beta<r\le\sigma-\theta$ such that $\beta$ is the only point of
$Z$ in $B(\sigma,r)$.
\end{enumerate}
\end{lemma}
\begin{proof}
The set $F$ of points of $Z$ in the closed disc of radius
$\sigma-\theta$ about $\sigma$ is finite. If $F$ is empty, case (1)
holds with $D=\sigma-\theta$ and $\rho_1=\rho_0$. Otherwise let $d$ be the
least distance from $\sigma$ to $F$; the disc $B(\sigma,d)$ meets no point
of $Z$. If some $\rho_1\in Z$ has $\sigma-\Re\rho_1<d$, case (1) holds
with $D=d$. If not, every point of $F$ at distance $d$ equals the real
point $\sigma-d$, so $\beta=\sigma-d\in Z$; moreover $d<\sigma-\theta$,
since $d=\sigma-\theta$ would contradict $\sigma-\Re\rho_0<\sigma-\theta$.
Taking $r$ between $d$ and the next distance (capped at $\sigma-\theta$)
gives case (2).
\end{proof}

\begin{lemma}[Isolating a nearest pole]\label{an:nearest}
\lean{Deduction.landau_nearest}
Let $Z$ be as in Lemma~\ref{an:dichotomy}, $\Re\rho<\sigma$ on $Z$,
$D<\sigma$ and $\rho_1\in Z$ with $\sigma-\Re\rho_1<D$. Then there are
$\gamma\in\R$, $\rho\in Z$ and $r\le D$ with $|\sigma+i\gamma-\rho|<r$ and
$\rho$ the only point of $Z$ in $B(\sigma+i\gamma,r)$.
\end{lemma}
\begin{proof}
Only finitely many points of $Z$ are relevant; for all but finitely many
$\gamma$ the point $\sigma+i\gamma$ is not equidistant from two of them,
and $\gamma$ near $\Im\rho_1$ brings a pole within distance $D$.
\end{proof}

\begin{lemma}[Coefficients near a simple pole]\label{an:pole-asym}
\lean{Deduction.landau_pole_asym, Deduction.landau_cauchy}
Let $\Phi$ be holomorphic on $B(z,r)\setminus\{\rho\}$, with $\rho\in
B(z,r)$, $\rho\ne z$, and $\Phi(s)=c/(s-\rho)+h(s)$ near $\rho$ with $h$
holomorphic at $\rho$. For $0<r'<r$ there is $M$ with
$|c_k(z)-c/(z-\rho)^{k+1}|\le M/r'^k$ for all $k$. If $\Phi$ is
holomorphic on all of $B(z,r)$, then $|\Phi^{(k)}(z)/k!|\le M/r'^k$ for all
$k$ (Cauchy's estimate).
\end{lemma}
\begin{proof}
Subtract the polar part, whose $k$-th coefficient at $z$ is
$c/(z-\rho)^{k+1}$, and apply Cauchy's estimate on the circle of radius
$r'$ to the holomorphic remainder.
\end{proof}

\begin{lemma}[Growth contradiction]\label{an:growth-contra}
\lean{Deduction.landau_growth_contra, Deduction.landau_window_contra}
Let $c>0$ and $0<d<e\le E$, and let $M,C\in\R$. It is impossible that for
every $\eta>0$ and $U_0$ there are $U\ge U_0$ and an integer $k\ge U$ with
$c/d^k\le M/e^k+Ce^{\eta U}/E^k$.
\end{lemma}
\begin{proof}
With $L=\log(e/d)>0$ and $\eta=L/2$, one has
$C'e^{\eta U}/e^k<c/d^k$ for all large $U$ and all $k\ge U$, where
$C'=|M|+|C|$; and $M/e^k+Ce^{\eta U}/E^k\le C'e^{\eta U}/e^k$.
\end{proof}

\begin{lemma}[Landau, window form for pole data]\label{an:landau-window}
\lean{Deduction.landau_window}
\uses{an:poledata, an:window-key, an:pole-re-le, an:dichotomy, an:nearest, an:pole-asym, an:growth-contra, an:windowhyp}
Assume $\mathrm{PoleData}(g,\Phi,Z)$, $0<\theta<1$, $\delta>0$ and
$\mathrm{WH}_U(g,\theta,\delta)$. Then $\Re\rho\le\theta$ for every
$\rho\in Z$.
\end{lemma}
\begin{proof}
\uses{an:window-key, an:pole-re-le, an:dichotomy, an:nearest, an:pole-asym, an:growth-contra}
Suppose $\rho_0\in Z$ has $\Re\rho_0>\theta$. Take $\sigma>2$ and $C$
from Lemma~\ref{an:window-key}; by Lemma~\ref{an:pole-re-le} all of $Z$
lies in $\Re s<\sigma$. Apply Lemma~\ref{an:dichotomy}.

In case (2), the residue $c$ at $\beta$ is real and positive, and
Lemma~\ref{an:pole-asym} at $z=\sigma$ with $r'=(\sigma-\beta+r)/2$ gives
$\Re c_k(\sigma)\ge c/(\sigma-\beta)^{k+1}-M/r'^k$. With $\gamma=0$,
Lemma~\ref{an:window-key} bounds $\Re c_k(\sigma)$ by
$Ce^{\eta U}/(\sigma-\theta)^k$ for one $k\ge U$ per window. Since
$\sigma-\beta<r'<r\le\sigma-\theta$, Lemma~\ref{an:growth-contra} gives a
contradiction.

In case (1), Lemma~\ref{an:nearest} gives $\gamma$ and a unique nearest
pole $\rho$ to $z=\sigma+i\gamma$ at distance $d<r\le D$. With
$r'=(d+r)/2$, Lemma~\ref{an:pole-asym} gives
$|c_k(z)|\ge|c|/d^{k+1}-M/r'^k$, and Cauchy's estimate on the pole-free
disc $B(\sigma,D)$ gives $\Re c_k(\sigma)\ge-M'/r'^k$. Adding and
comparing with Lemma~\ref{an:window-key} again contradicts
Lemma~\ref{an:growth-contra}, as $d<r'\le D\le\sigma-\theta$.

The lower bounds hold for every $k$, so one admissible $k$ per window is
enough.
\end{proof}

\begin{lemma}\label{an:lemmaL}
\lean{Deduction.lemmaL}
\uses{an:charfamily, an:windowhyp, an:poledata-lemma, an:landau-window, an:windowhypu, an:nonvanish-re1}
Let $F$ be a character family, $0<\theta<1$ and $\delta>0$. If
$\mathrm{WH}(F,\theta,\delta)$, then $F$ is zero-free above $\theta$.
\end{lemma}
\begin{proof}
\uses{an:poledata-lemma, an:landau-window, an:windowhypu, an:nonvanish-re1}
Let $L(s,\chi_i)=0$ with $\Re s>\theta$. If $\Re s\ge1$ this contradicts
Lemma~\ref{an:nonvanish-re1}. Otherwise $s\in Z(F)$, and
Lemma~\ref{an:landau-window}, applied to the pole data of
Lemma~\ref{an:poledata-lemma} and the hypothesis transported by
Lemma~\ref{an:windowhypu}, gives $\Re s\le\theta$, a contradiction.
\end{proof}

\begin{theorem}[Windowed Landau theorem]\label{an:rpow-window}
\lean{Deduction.zeroFree_of_rpow_window}
\uses{an:charfamily, an:lemmaL, an:windowhyp-rpow}
Let $F$ be a character family, $0<\theta<1$, $\delta>0$ and $C\in\R$.
Suppose that for every $Y_0$ there is $Y\ge Y_0$ with
$G_F(y)\le Cy^{\theta}$ for all $y\in[Y,Y^{1+\delta}]$. Then
$L(s,\chi_i)\ne0$ for every $i$ and every $s$ with $\Re s>\theta$.
\end{theorem}
\begin{proof}
\uses{an:lemmaL, an:windowhyp-rpow}
Lemma~\ref{an:windowhyp-rpow} followed by Lemma~\ref{an:lemmaL}.
\end{proof}

\begin{remark}[why positivity matters]\label{an:rem-positivity}
The theorem uses only an upper bound for $G_F$. The proof needs the
residue of the transform at a real zero to be positive
(Definition~\ref{an:poledata}(4)), and this comes from $a_i>0$ in
Lemma~\ref{an:poledata-lemma}. This is why the archive must produce a
combination of $\Gz$ and the $\Re g_{\psi}$ with positive weights.
\end{remark}

\section{The window criterion for every level}

\begin{lemma}[A level as a character family]\label{an:level-family}
\lean{UpperTheory.Comparison₂.level_family, UpperTheory.Comparison₂.gChi_eq}
\uses{an:level, an:charfamily, an:gz-eq, an:zerofree}
For every level $\Lambda$ there is a character family $F$ with
$G_F=G_{\Lambda}$ and, for every $\theta$, $F$ zero-free above $\theta$ if
and only if $\Lambda$ is zero-free above $\theta$.
\end{lemma}
\begin{proof}
\uses{an:gz-eq, an:level, an:charfamily}
Index $F$ by $\Kept(\Lambda)$, with the character $\psi$ itself, weight
$1$ for $\psi=\mathbf 1$ and $\operatorname{wt}(\psi)$ otherwise. All
weights are positive. Admissibility holds because a primitive $\psi$
with $\psi=1$ has modulus $1$. The two versions of the smoothed sums
agree: $g_{\psi}=\mathcal G_{\psi}$ for primitive $\psi$, and $\Gz=G$
(Lemma~\ref{an:gz-eq}), so splitting off the term $\mathbf 1$ gives
$G_F=G_{\Lambda}$.
\end{proof}

\begin{theorem}[Every level satisfies the window criterion]\label{an:formb-thm}
\lean{UpperTheory.Comparison₂.formB}
\uses{an:formb, an:level-family, an:rpow-window}
Every level $\Lambda$ satisfies the window criterion at every $\lambda>1$.
\end{theorem}
\begin{proof}
\uses{an:level-family, an:rpow-window}
Given $\lambda>1$, $\theta\in(0,1)$ and windows of exponent $\lambda$ for
$\{G_{\Lambda}(x)\le x^{\theta}\}$, pass to the family $F$ of
Lemma~\ref{an:level-family} and apply Theorem~\ref{an:rpow-window} with
$\delta=\lambda-1$ and $C=1$.
\end{proof}

\section{From the Chebyshev inequality to zero-freeness}



\begin{remark}
The exponents $e_D,e_L$ are chosen before $\varepsilon$ in the second
statement, so the exponent $1/2+(e_D+e_L+7)\varepsilon$ can be pushed to
$1/2$. The same family and the same relation $\mathrm{Tied}$ appear in
both statements, so the first applies to the family produced by the
second.
\end{remark}

\begin{lemma}\label{an:small-eps}
\lean{Formalism.HC.exists_small_eps}
For $C\in\R$, $\delta>0$ and $\varepsilon_1>0$ there is
$\varepsilon\in(0,\varepsilon_1)$ with $C\varepsilon\le\delta$.
\end{lemma}
\begin{proof}
Take $\varepsilon=\min(\varepsilon_1/2,\delta/(|C|+1))$.
\end{proof}

\begin{lemma}[A level from tied windows]\label{an:level-tiedwin}
\lean{Formalism.HC.level_of_tiedWin}
\uses{geoa:efamc, an:formb, an:small-eps, an:windows-mono, an:zerofree}
Let $\Lambda$ be a level and $\lambda\in\R$. Assume the Chebyshev
inequality on windows for $\Lambda$, the window criterion for $\Lambda$ at
$\lambda$, and the Chebyshev inequality at window scale $\lambda$ for
$\Lambda$. Then $\mathrm{ZeroFreeLevel}(\Lambda)$.
\end{lemma}
\begin{proof}
\uses{an:small-eps, an:windows-mono}
Let $\psi\in\Kept(\Lambda)$ and $\Re s>1/2$. Put
$m=\min(\Re s-1/2,1/2)>0$ and $\theta=1/2+m/2\in(0,1)$, so
$\theta<\Re s$. By Lemma~\ref{an:small-eps} choose $\varepsilon$ below
$\varepsilon_1$ with $(e_D+e_L+7)\varepsilon\le m/2$, and take the family
given by the window-scale inequality. The inequality on windows gives
$G_{\Lambda}(x)\le x^{1/2+(e_D+e_L+7)\varepsilon}\le x^{\theta}$ for large
$x\ge1$ in $\mathrm{Win}$; by Lemma~\ref{an:windows-mono}, the set
$\{G_{\Lambda}(x)\le x^{\theta}\}$ inherits the windows of $\mathrm{Win}$.
The window criterion gives $L(s,\psi)\ne0$.
\end{proof}

\begin{proposition}\label{an:zerofree-tiedwin}
\lean{Formalism.HC.zeroFreeOn_of_tiedWin, Formalism.HC.zeroFreeOn_of_tiedWin_cyc}
\uses{an:level-tiedwin, an:zerofree-levels, an:level-cyc, an:formb}
Let $Q$ be a property of primitive characters and $\lambda>1$. Suppose
that for every $\psi$ with $Q(\psi)$ the cyclic level $\Lambda_{\psi}$
satisfies the Chebyshev inequality on windows, the window criterion, and
the Chebyshev inequality at window scale $\lambda$. Then
$\mathrm{ZeroFreeOn}(Q)$. (The general form allows any level containing
$\psi$ and any $\lambda>1$ depending on $\psi$.)
\end{proposition}
\begin{proof}
\uses{an:level-tiedwin, an:zerofree-levels, an:level-cyc}
Lemma~\ref{an:level-tiedwin} gives $\mathrm{ZeroFreeLevel}(\Lambda_{\psi})$,
using the window criterion at $\lambda$; since
$\psi\in\Kept(\Lambda_{\psi})$, Lemma~\ref{an:zerofree-levels} concludes.
\end{proof}



\begin{remark}
$\mathrm{TopG}$ is quantified over every host requirement $G$, while the
chain uses it only for $G=G_{\chi}(\psi)$, the requirement that the host
satisfy $\mathrm{EAtH}$ for $\Lambda_{\psi}$ with constant $4$ together with
a bound $\mathrm{VVle}$ on Tate--Shafarevich lengths. Its content lies in
the chapters on the Chebyshev inequality; this chapter uses it only
through its statement.
\end{remark}

\begin{lemma}\label{an:cls}
\lean{UpperTheory.Top.hypV_of_cls, UpperTheory.Top.eFamC_of_cls}
% uses-lean: UpperTheory.Top.GChi, UpperTheory.Top.FReq.cls, UpperTheory.Top.HypV, Formalism.MainNotions.FamP
\uses{geoa:famcore, host:gchi} % resolved
\uses{geoa:efamc, an:level-cyc}
Let $\psi$ be a primitive character, $\lambda\in\R$, and $\mathcal F$ a
family in $\operatorname{cls}_{\lambda}(G_{\chi}(\psi))$. Then $\mathcal F$
satisfies $\mathrm{HypV}$ and $\mathrm{EFamC}(\Lambda_{\psi},\mathcal F)$
(with constant $4$).
\end{lemma}
\begin{proof}
Both are components of $G_{\chi}(\psi)$ at each datum beyond the family
threshold, after identifying the datum's scale and parameter with $x$ and
$\varepsilon$.
\end{proof}

\begin{theorem}[Zero-freeness for all primitive characters]\label{an:zerofreeall-topg}
\lean{UpperTheory.Top.zeroFreeAll_of_topG}
% uses-lean: UpperTheory.Comparison₂.level_gm_le_on_family_scales, GlobalTheory.EX.HostExW, UpperTheory.Top.GChi, UpperTheory.Top.FReq.req, Formalism.MainNotions.ChainTied
\uses{geoa:tied, geob:levelgm, host:gchi, host:host-exw} % resolved
\uses{geob:topg, an:zerofree-tiedwin, an:formb-thm, an:cls, an:zerofree}
Let $\lambda>1$. Assume $\mathrm{TopG}$, and assume that for every
primitive character $\psi$ hosts satisfying $G_{\chi}(\psi)$ exist on windows
of exponent $\lambda$. Then $\mathrm{ZeroFreeAll}$: $L(s,\psi)\ne0$ for
every primitive Dirichlet character $\psi$ and every $s$ with $\Re s>1/2$.
\end{theorem}
\begin{proof}
\uses{an:zerofree-tiedwin, an:formb-thm, an:cls, geob:topg}
Apply Proposition~\ref{an:zerofree-tiedwin} with $Q$ true, families of the
parametrised type $\mathrm{FamP}$ and $\mathrm{Tied}=\mathrm{ChainTied}$.
Its three hypotheses for $\Lambda_{\psi}$ are: the Chebyshev inequality on
windows, which is the archive's
\texttt{level\_gm\_le\_on\_family\_scales} (proved for every level in the
chapters on the Chebyshev inequality); the window criterion,
Theorem~\ref{an:formb-thm}; and the inequality at window scale
$\lambda$, obtained from $\mathrm{TopG}$ at $G=G_{\chi}(\psi)$, whose
hypothesis $\mathrm{HypV}$ is supplied by Lemma~\ref{an:cls}, and whose
output families lie in $\operatorname{cls}_{\lambda}(G_{\chi}(\psi))$, so
they have windows of exponent $\lambda$ and satisfy $\mathrm{EFamC}$ by
Lemma~\ref{an:cls}.
\end{proof}

\begin{remark}
The archive applies this theorem with $\lambda=3/2$. Passing from
$\mathrm{ZeroFreeAll}$ (a statement about primitive characters) to all
Dirichlet characters and to the Riemann hypothesis is done elsewhere in
the archive.
\end{remark}

\section{Parallel and unused developments}

\begin{definition}[Non-windowed criterion]\label{an:chebgrowth}
\lean{Formalism.HC.Level.ChebyshevGrowthImpliesZeroFree}
\uses{an:level, an:zerofree}
For a level $\Lambda$, $\mathrm{ChebyshevGrowthImpliesZeroFree}(\Lambda)$
says: for every $\theta\in(0,1)$ and $y_0$, if $G_{\Lambda}(y)\le y^{\theta}$
for all $y\ge y_0$, then $\Lambda$ is zero-free above $\theta$.
\end{definition}

\begin{lemma}\label{an:chebgrowth-lemma}
\lean{Formalism.HC.Landau.zero_free_of_Chebyshev_growth, Formalism.HC.Landau.landau_family}
\uses{an:chebgrowth, an:re-gchi}
Every level satisfies $\mathrm{ChebyshevGrowthImpliesZeroFree}$.
\end{lemma}
\begin{proof}
\uses{an:re-gchi}
Multiply by $3o(\Lambda)+1$ to obtain integer multiplicities ($3o+1$ for
$\psi_K$, $1$ for the other nontrivial kept characters) and apply a
Landau theorem for $R\,\Gz+\sum_{\xi} G^{\R}_{\chi_{\xi}}$ over a finite
inversion-closed family of nontrivial characters, proved in
\arch{Formalism/SharpRestriction/RamificationFrobenius/TotientGaloisClasses/Landau.lean}
from the product $\zeta_1^R\prod L(\cdot,\chi_{\xi})$; Lemma~\ref{an:re-gchi}
converts $\Re g$ to $G^{\R}$.
\end{proof}

\begin{remark}[a file named after Landau]\label{an:rem-landaufile}
The file \arch{.../TotientGaloisClasses/Landau.lean} proves only the
non-windowed Lemma~\ref{an:chebgrowth-lemma}. It feeds the alternative
route \texttt{level\_of\_tiedAll} / \texttt{zeroFreeOn\_of\_tiedAll},
which needs the tie on all large scales. The chain used by
Theorem~\ref{an:zerofreeall-topg} goes through Theorem~\ref{an:formb-thm}
and the separate development in \arch{.../DedekindLFunctions/LandauTheorems/}.
\end{remark}

\begin{remark}[definitions without uses]\label{an:rem-unused}
\arch{Formalism/StatementOC.lean} also defines statements
\texttt{OCStep}, \texttt{OCStepUpper}, \texttt{OCAvg}, \texttt{OCAvgUpper}
and \texttt{RiemannHypothesisOfLowerBounds} about one-sided bounds for
$\Gz-G^{\R}_{\chi}$ with real characters; none of them is referenced
elsewhere in the archive. The same holds for \texttt{Level.LandauWin}
(\arch{.../TotientGaloisClasses/Window.lean}). The predicates
\texttt{LandauZeroFreeCriterion}, \texttt{GmLevel} and \texttt{famD} in
\arch{.../DedekindLFunctions/LandauZeroFreeCriterion.lean} concern even
characters modulo $2^{j+2}$ and are not used by the chain above.
\end{remark}
